
\documentclass{article}
\usepackage{colortbl}
\usepackage{makecell}
\usepackage{multirow}
\usepackage{supertabular}

\begin{document}

\newcounter{utterance}

\centering \large Interaction Transcript for game `codenames', experiment `opponent\_difficult', episode 1 with deepseek{-}v4{-}pro.
\vspace{24pt}

{ \footnotesize  \setcounter{utterance}{1}
\setlength{\tabcolsep}{0pt}
\begin{supertabular}{c@{$\;$}|p{.15\linewidth}@{}p{.15\linewidth}p{.15\linewidth}p{.15\linewidth}p{.15\linewidth}p{.15\linewidth}}
    \# & $\;$A & \multicolumn{4}{c}{Game Master} & $\;\:$B\\
    \hline

    \theutterance \stepcounter{utterance}  
    & & & \multicolumn{2}{p{0.3\linewidth}}{
        \cellcolor[rgb]{0.9,0.9,0.9}{
            \makecell[{{p{\linewidth}}}]{
                \texttt{\tiny{[GM$|$GM]}}
                {'hidden': {'team': ['nu', 'pārsegs', 'savienojums', 'ziloņkauls', 'atlantīda', 'galva', 'zāle', 'pātagu', 'saturns'], 'innocent': ['ceļojums', 'amazone', 'bērns', 'kaķis', 'automašīna', 'lāpa', 'telpa'], 'opponent': ['sniegs', 'pudele', 'kļava', 'dēlis', 'vēstniecība', 'antarktīda', 'aste', 'mārciņa'], 'assassin': ['nāve']}, 'revealed': {'team': {'team': [], 'innocent': [], 'opponent': [], 'assassin': []}, 'opponent': {'team': [], 'innocent': [], 'opponent': [], 'assassin': []}}}
            }
        }
    }
    & & \\ \\

    \theutterance \stepcounter{utterance}  
    & & \multicolumn{4}{p{0.6\linewidth}}{
        \cellcolor[rgb]{0.9,0.9,0.9}{
            \makecell[{{p{\linewidth}}}]{
                \texttt{\tiny{[P1$\langle$GM]}}
                \texttt{Sniedziet vienu vienvārdu norādi, kas ir saistīta ar kādu no jūsu komandas vārdiem no tālāk dotā saraksta. Jūs varat izvēlēties, uz cik un kādiem vārdiem attiecas jūsu norāde.} \\
\texttt{Mērķis ir panākt, lai jūsu partneris uzmin visus jūsu komandas vārdus, pamatojoties uz jūsu dotajām norādēm, pirms otra komanda ir atklājusi visus savus vārdus. Ja jūsu partneris uzmin "pretinieka" vai "novirzītāja" vārdu, viņa uzminēšanas kārta beidzas, un pretinieku komanda var veikt savu kārtu. Ja jūsu partneris uzmin "nogalināšanas" vārdu, spēle beidzas un jūsu komanda uzreiz zaudē.} \\
\texttt{Jo vairāk vārdu, uz kuriem attiecas jūsu norāde, jo ātrāk jūs varat uzvarēt spēli, bet jo grūtāk asociācijas var būt uzminēt jūsu partnerim, tāpēc izvēlieties savu norādes vārdu gudri. Norādes vārdam ir jābūt semantiski saistītam ar mērķa vārdiem, tas nedrīkst būt viens no sarakstos iekļautajiem vārdiem vai saturēt to daļas.} \\
\\ 
\texttt{Vienvārdu norādi un ar komatu atdalīto saistīto mērķa vārdu sarakstu vienmēr sniedziet šādā formātā un atbildi veidojiet pēc iespējas īsāku, nekad neiekļaujiet citu tekstu, kas nav prasīts šajā formā:} \\
\\ 
\texttt{PADOMS: <VĀRDS>} \\
\texttt{MĒRĶI: <MĒRĶI>} \\
\\ 
\texttt{Jūsu komandas vārdi ir: nu, pārsegs, savienojums, ziloņkauls, atlantīda, galva, zāle, pātagu, saturns.} \\
\texttt{Jūsu pretinieka komandas vārdi ir: sniegs, pudele, kļava, dēlis, vēstniecība, antarktīda, aste, mārciņa.} \\
\texttt{Novirzītāja vārdi ir: ceļojums, amazone, bērns, kaķis, automašīna, lāpa, telpa.} \\
\texttt{Nogalināšanas vārdi ir: nāve.} \\
            }
        }
    }
    & \\ \\

    \theutterance \stepcounter{utterance}  
    & \multicolumn{4}{p{0.6\linewidth}}{
        \cellcolor[rgb]{0.9,0.9,0.9}{
            \makecell[{{p{\linewidth}}}]{
                \texttt{\tiny{[P1$\rangle$GM]}}
                \texttt{PADOMS: grieķu} \\
\texttt{MĒRĶI: atlantīda, saturns} \\
            }
        }
    }
    & & \\ \\

    \theutterance \stepcounter{utterance}  
    & & & \multicolumn{2}{p{0.3\linewidth}}{
        \cellcolor[rgb]{0.9,0.9,0.9}{
            \makecell[{{p{\linewidth}}}]{
                \texttt{\tiny{[GM$|$GM]}}
                \texttt{grieķu} \\
            }
        }
    }
    & & \\ \\

    \theutterance \stepcounter{utterance}  
    & & & \multicolumn{2}{p{0.3\linewidth}}{
        \cellcolor[rgb]{0.9,0.9,0.9}{
            \makecell[{{p{\linewidth}}}]{
                \texttt{\tiny{[GM$|$GM]}}
                ['atlantīda', 'saturns']
            }
        }
    }
    & & \\ \\

    \theutterance \stepcounter{utterance}  
    & & & \multicolumn{2}{p{0.3\linewidth}}{
        \cellcolor[rgb]{0.9,0.9,0.9}{
            \makecell[{{p{\linewidth}}}]{
                \texttt{\tiny{[GM$|$GM]}}
                {'word': 'atlantīda', 'assignment': 'team'}
            }
        }
    }
    & & \\ \\

    \theutterance \stepcounter{utterance}  
    & & & \multicolumn{2}{p{0.3\linewidth}}{
        \cellcolor[rgb]{0.9,0.9,0.9}{
            \makecell[{{p{\linewidth}}}]{
                \texttt{\tiny{[GM$|$GM]}}
                {'word': 'saturns', 'assignment': 'team'}
            }
        }
    }
    & & \\ \\

    \theutterance \stepcounter{utterance}  
    & & \multicolumn{4}{p{0.6\linewidth}}{
        \cellcolor[rgb]{0.9,0.9,0.9}{
            \makecell[{{p{\linewidth}}}]{
                \texttt{\tiny{[GM$\rangle$P2]}}
                \texttt{Sniedziet ar komatiem atdalītu sarakstu ar līdz 2 vārdiem no turpmāk minētā saraksta, kas vislabāk attiecas uz vārdu 'grieķu' vai ir visciešāk saistīti ar to. Vienmēr sāciet minējumu sarakstu ar MINĒJUMS:  un atbildē neiekļaujiet nekādu citu tekstu.} \\
\\ 
\texttt{lāpa, pudele, antarktīda, aste, mārciņa, pārsegs, vēstniecība, ziloņkauls, atlantīda, kļava, amazone, bērns, savienojums, zāle, ceļojums, telpa, kaķis, pātagu, sniegs, nāve, automašīna, dēlis, galva, saturns, nu} \\
            }
        }
    }
    & \\ \\

    \theutterance \stepcounter{utterance}  
    & & & \multicolumn{4}{p{0.6\linewidth}}{
        \cellcolor[rgb]{0.9,0.9,0.9}{
            \makecell[{{p{\linewidth}}}]{
                \texttt{\tiny{[GM$\langle$P2]}}
                \texttt{MINĒJUMS: atlantīda, amazone} \\
            }
        }
    }
     \\ \\

    \theutterance \stepcounter{utterance}  
    & & & \multicolumn{2}{p{0.3\linewidth}}{
        \cellcolor[rgb]{0.9,0.9,0.9}{
            \makecell[{{p{\linewidth}}}]{
                \texttt{\tiny{[GM$|$GM]}}
                ['atlantīda', 'amazone']
            }
        }
    }
    & & \\ \\

    \theutterance \stepcounter{utterance}  
    & & & \multicolumn{2}{p{0.3\linewidth}}{
        \cellcolor[rgb]{0.9,0.9,0.9}{
            \makecell[{{p{\linewidth}}}]{
                \texttt{\tiny{[GM$|$GM]}}
                {'word': 'atlantīda', 'assignment': 'team'}
            }
        }
    }
    & & \\ \\

    \theutterance \stepcounter{utterance}  
    & & & \multicolumn{2}{p{0.3\linewidth}}{
        \cellcolor[rgb]{0.9,0.9,0.9}{
            \makecell[{{p{\linewidth}}}]{
                \texttt{\tiny{[GM$|$GM]}}
                {'word': 'atlantīda', 'assignment': 'team'}
            }
        }
    }
    & & \\ \\

    \theutterance \stepcounter{utterance}  
    & & & \multicolumn{2}{p{0.3\linewidth}}{
        \cellcolor[rgb]{0.9,0.9,0.9}{
            \makecell[{{p{\linewidth}}}]{
                \texttt{\tiny{[GM$|$GM]}}
                {'word': 'amazone', 'assignment': 'innocent'}
            }
        }
    }
    & & \\ \\

    \theutterance \stepcounter{utterance}  
    & & & \multicolumn{2}{p{0.3\linewidth}}{
        \cellcolor[rgb]{0.9,0.9,0.9}{
            \makecell[{{p{\linewidth}}}]{
                \texttt{\tiny{[GM$|$GM]}}
                \texttt{amazone} \\
            }
        }
    }
    & & \\ \\

    \theutterance \stepcounter{utterance}  
    & & & \multicolumn{2}{p{0.3\linewidth}}{
        \cellcolor[rgb]{0.9,0.9,0.9}{
            \makecell[{{p{\linewidth}}}]{
                \texttt{\tiny{[GM$|$GM]}}
                {'word': 'pudele', 'assignment': 'opponent'}
            }
        }
    }
    & & \\ \\

    \theutterance \stepcounter{utterance}  
    & & & \multicolumn{2}{p{0.3\linewidth}}{
        \cellcolor[rgb]{0.9,0.9,0.9}{
            \makecell[{{p{\linewidth}}}]{
                \texttt{\tiny{[GM$|$GM]}}
                {'word': 'dēlis', 'assignment': 'opponent'}
            }
        }
    }
    & & \\ \\

    \theutterance \stepcounter{utterance}  
    & & & \multicolumn{2}{p{0.3\linewidth}}{
        \cellcolor[rgb]{0.9,0.9,0.9}{
            \makecell[{{p{\linewidth}}}]{
                \texttt{\tiny{[GM$|$GM]}}
                {'hidden': {'team': ['nu', 'pārsegs', 'savienojums', 'ziloņkauls', 'galva', 'zāle', 'pātagu', 'saturns'], 'innocent': ['ceļojums', 'bērns', 'kaķis', 'automašīna', 'lāpa', 'telpa'], 'opponent': ['sniegs', 'kļava', 'vēstniecība', 'antarktīda', 'aste', 'mārciņa'], 'assassin': ['nāve']}, 'revealed': {'team': {'team': ['atlantīda'], 'innocent': ['amazone'], 'opponent': [], 'assassin': []}, 'opponent': {'team': [], 'innocent': [], 'opponent': ['pudele', 'dēlis'], 'assassin': []}}}
            }
        }
    }
    & & \\ \\

    \theutterance \stepcounter{utterance}  
    & & \multicolumn{4}{p{0.6\linewidth}}{
        \cellcolor[rgb]{0.9,0.9,0.9}{
            \makecell[{{p{\linewidth}}}]{
                \texttt{\tiny{[P1$\langle$GM]}}
                \texttt{Vārds atlantīda tika uzminēts pareizi. Vārds amazone tika uzminēts, bet tas ir nevainojams vārds. Jūsu komandas biedra gājiens beidzās tur.} \\
\texttt{Tagad sniedziet vēl vienu norādi, kas attiecas uz dažiem no jūsu atlikušajiem komandas vārdiem, un saistīto mērķa vārdu sarakstu. Neaizmirstiet sākt savu norādi ar 'PADOMS: ', pēc tam ievietot jaunu rindu un sākt ar komatiem atdalīto mērķa vārdu sarakstu ar 'MĒRĶI: '. Piezīme: salīdzinājumā ar iepriekšējiem pieprasījumiem daži vārdi no sarakstiem ir izņemti.} \\
\\ 
\texttt{Jūsu atlikušie komandas vārdi ir: nu, pārsegs, savienojums, ziloņkauls, galva, zāle, pātagu, saturns.} \\
\texttt{Atlikušie pretinieka vārdi ir: sniegs, pudele, kļava, dēlis, vēstniecība, antarktīda, aste, mārciņa.} \\
\texttt{Atlikušie traucētājvārdi ir: ceļojums, bērns, kaķis, automašīna, lāpa, telpa.} \\
\texttt{Atlikušie slepkavības vārdi ir: nāve.} \\
            }
        }
    }
    & \\ \\

    \theutterance \stepcounter{utterance}  
    & \multicolumn{4}{p{0.6\linewidth}}{
        \cellcolor[rgb]{0.9,0.9,0.9}{
            \makecell[{{p{\linewidth}}}]{
                \texttt{\tiny{[P1$\rangle$GM]}}
                \texttt{PADOMS: kosmoss} \\
\texttt{MĒRĶI: saturns, savienojums} \\
            }
        }
    }
    & & \\ \\

    \theutterance \stepcounter{utterance}  
    & & & \multicolumn{2}{p{0.3\linewidth}}{
        \cellcolor[rgb]{0.9,0.9,0.9}{
            \makecell[{{p{\linewidth}}}]{
                \texttt{\tiny{[GM$|$GM]}}
                \texttt{kosmoss} \\
            }
        }
    }
    & & \\ \\

    \theutterance \stepcounter{utterance}  
    & & & \multicolumn{2}{p{0.3\linewidth}}{
        \cellcolor[rgb]{0.9,0.9,0.9}{
            \makecell[{{p{\linewidth}}}]{
                \texttt{\tiny{[GM$|$GM]}}
                ['saturns', 'savienojums']
            }
        }
    }
    & & \\ \\

    \theutterance \stepcounter{utterance}  
    & & & \multicolumn{2}{p{0.3\linewidth}}{
        \cellcolor[rgb]{0.9,0.9,0.9}{
            \makecell[{{p{\linewidth}}}]{
                \texttt{\tiny{[GM$|$GM]}}
                {'word': 'saturns', 'assignment': 'team'}
            }
        }
    }
    & & \\ \\

    \theutterance \stepcounter{utterance}  
    & & & \multicolumn{2}{p{0.3\linewidth}}{
        \cellcolor[rgb]{0.9,0.9,0.9}{
            \makecell[{{p{\linewidth}}}]{
                \texttt{\tiny{[GM$|$GM]}}
                {'word': 'savienojums', 'assignment': 'team'}
            }
        }
    }
    & & \\ \\

    \theutterance \stepcounter{utterance}  
    & & \multicolumn{4}{p{0.6\linewidth}}{
        \cellcolor[rgb]{0.9,0.9,0.9}{
            \makecell[{{p{\linewidth}}}]{
                \texttt{\tiny{[GM$\rangle$P2]}}
                \texttt{Vārds atlantīda tika uzminēts pareizi. Vārds amazone tika uzminēts, bet tas ir nevainojams vārds. Jūsu gājiens beidzās tur.} \\
\texttt{Tagad sniedziet vēl vienu ar komatiem atdalītu sarakstu, kurā ir vismaz 1 un ne vairāk kā 2 vārdi no turpmāk minētā vārdu saraksta, kas vislabāk sasaucas vai ir visciešāk saistīti ar vārdu 'kosmoss'. Neaizmirstiet atbildi sākt ar MINĒJUMS: . Piezīme: salīdzinājumā ar iepriekšējiem pieprasījumiem daži vārdi no saraksta ir izņemti.} \\
\\ 
\texttt{lāpa, antarktīda, aste, mārciņa, pārsegs, vēstniecība, ziloņkauls, kļava, bērns, savienojums, zāle, ceļojums, telpa, kaķis, pātagu, sniegs, nāve, automašīna, galva, saturns, nu} \\
            }
        }
    }
    & \\ \\

    \theutterance \stepcounter{utterance}  
    & & & \multicolumn{4}{p{0.6\linewidth}}{
        \cellcolor[rgb]{0.9,0.9,0.9}{
            \makecell[{{p{\linewidth}}}]{
                \texttt{\tiny{[GM$\langle$P2]}}
                \texttt{MINĒJUMS: saturns, telpa} \\
            }
        }
    }
     \\ \\

    \theutterance \stepcounter{utterance}  
    & & & \multicolumn{2}{p{0.3\linewidth}}{
        \cellcolor[rgb]{0.9,0.9,0.9}{
            \makecell[{{p{\linewidth}}}]{
                \texttt{\tiny{[GM$|$GM]}}
                ['saturns', 'telpa']
            }
        }
    }
    & & \\ \\

    \theutterance \stepcounter{utterance}  
    & & & \multicolumn{2}{p{0.3\linewidth}}{
        \cellcolor[rgb]{0.9,0.9,0.9}{
            \makecell[{{p{\linewidth}}}]{
                \texttt{\tiny{[GM$|$GM]}}
                {'word': 'saturns', 'assignment': 'team'}
            }
        }
    }
    & & \\ \\

    \theutterance \stepcounter{utterance}  
    & & & \multicolumn{2}{p{0.3\linewidth}}{
        \cellcolor[rgb]{0.9,0.9,0.9}{
            \makecell[{{p{\linewidth}}}]{
                \texttt{\tiny{[GM$|$GM]}}
                {'word': 'saturns', 'assignment': 'team'}
            }
        }
    }
    & & \\ \\

    \theutterance \stepcounter{utterance}  
    & & & \multicolumn{2}{p{0.3\linewidth}}{
        \cellcolor[rgb]{0.9,0.9,0.9}{
            \makecell[{{p{\linewidth}}}]{
                \texttt{\tiny{[GM$|$GM]}}
                {'word': 'telpa', 'assignment': 'innocent'}
            }
        }
    }
    & & \\ \\

    \theutterance \stepcounter{utterance}  
    & & & \multicolumn{2}{p{0.3\linewidth}}{
        \cellcolor[rgb]{0.9,0.9,0.9}{
            \makecell[{{p{\linewidth}}}]{
                \texttt{\tiny{[GM$|$GM]}}
                \texttt{telpa} \\
            }
        }
    }
    & & \\ \\

    \theutterance \stepcounter{utterance}  
    & & & \multicolumn{2}{p{0.3\linewidth}}{
        \cellcolor[rgb]{0.9,0.9,0.9}{
            \makecell[{{p{\linewidth}}}]{
                \texttt{\tiny{[GM$|$GM]}}
                {'word': 'antarktīda', 'assignment': 'opponent'}
            }
        }
    }
    & & \\ \\

    \theutterance \stepcounter{utterance}  
    & & & \multicolumn{2}{p{0.3\linewidth}}{
        \cellcolor[rgb]{0.9,0.9,0.9}{
            \makecell[{{p{\linewidth}}}]{
                \texttt{\tiny{[GM$|$GM]}}
                {'word': 'aste', 'assignment': 'opponent'}
            }
        }
    }
    & & \\ \\

    \theutterance \stepcounter{utterance}  
    & & & \multicolumn{2}{p{0.3\linewidth}}{
        \cellcolor[rgb]{0.9,0.9,0.9}{
            \makecell[{{p{\linewidth}}}]{
                \texttt{\tiny{[GM$|$GM]}}
                {'hidden': {'team': ['nu', 'pārsegs', 'savienojums', 'ziloņkauls', 'galva', 'zāle', 'pātagu'], 'innocent': ['ceļojums', 'bērns', 'kaķis', 'automašīna', 'lāpa'], 'opponent': ['sniegs', 'kļava', 'vēstniecība', 'mārciņa'], 'assassin': ['nāve']}, 'revealed': {'team': {'team': ['atlantīda', 'saturns'], 'innocent': ['amazone', 'telpa'], 'opponent': [], 'assassin': []}, 'opponent': {'team': [], 'innocent': [], 'opponent': ['pudele', 'dēlis', 'antarktīda', 'aste'], 'assassin': []}}}
            }
        }
    }
    & & \\ \\

    \theutterance \stepcounter{utterance}  
    & & \multicolumn{4}{p{0.6\linewidth}}{
        \cellcolor[rgb]{0.9,0.9,0.9}{
            \makecell[{{p{\linewidth}}}]{
                \texttt{\tiny{[P1$\langle$GM]}}
                \texttt{Vārds saturns tika uzminēts pareizi. Vārds telpa tika uzminēts, bet tas ir nevainojams vārds. Jūsu komandas biedra gājiens beidzās tur.} \\
\texttt{Tagad sniedziet vēl vienu norādi, kas attiecas uz dažiem no jūsu atlikušajiem komandas vārdiem, un saistīto mērķa vārdu sarakstu. Neaizmirstiet sākt savu norādi ar 'PADOMS: ', pēc tam ievietot jaunu rindu un sākt ar komatiem atdalīto mērķa vārdu sarakstu ar 'MĒRĶI: '. Piezīme: salīdzinājumā ar iepriekšējiem pieprasījumiem daži vārdi no sarakstiem ir izņemti.} \\
\\ 
\texttt{Jūsu atlikušie komandas vārdi ir: nu, pārsegs, savienojums, ziloņkauls, galva, zāle, pātagu.} \\
\texttt{Atlikušie pretinieka vārdi ir: sniegs, kļava, vēstniecība, antarktīda, aste, mārciņa.} \\
\texttt{Atlikušie traucētājvārdi ir: ceļojums, bērns, kaķis, automašīna, lāpa.} \\
\texttt{Atlikušie slepkavības vārdi ir: nāve.} \\
            }
        }
    }
    & \\ \\

    \theutterance \stepcounter{utterance}  
    & \multicolumn{4}{p{0.6\linewidth}}{
        \cellcolor[rgb]{0.9,0.9,0.9}{
            \makecell[{{p{\linewidth}}}]{
                \texttt{\tiny{[P1$\rangle$GM]}}
                \texttt{PADOMS: zilonis} \\
\texttt{MĒRĶI: ziloņkauls, galva} \\
            }
        }
    }
    & & \\ \\

    \theutterance \stepcounter{utterance}  
    & & & \multicolumn{2}{p{0.3\linewidth}}{
        \cellcolor[rgb]{0.9,0.9,0.9}{
            \makecell[{{p{\linewidth}}}]{
                \texttt{\tiny{[GM$|$GM]}}
                \texttt{zilonis} \\
            }
        }
    }
    & & \\ \\

    \theutterance \stepcounter{utterance}  
    & & & \multicolumn{2}{p{0.3\linewidth}}{
        \cellcolor[rgb]{0.9,0.9,0.9}{
            \makecell[{{p{\linewidth}}}]{
                \texttt{\tiny{[GM$|$GM]}}
                ['ziloņkauls', 'galva']
            }
        }
    }
    & & \\ \\

    \theutterance \stepcounter{utterance}  
    & & & \multicolumn{2}{p{0.3\linewidth}}{
        \cellcolor[rgb]{0.9,0.9,0.9}{
            \makecell[{{p{\linewidth}}}]{
                \texttt{\tiny{[GM$|$GM]}}
                {'word': 'ziloņkauls', 'assignment': 'team'}
            }
        }
    }
    & & \\ \\

    \theutterance \stepcounter{utterance}  
    & & & \multicolumn{2}{p{0.3\linewidth}}{
        \cellcolor[rgb]{0.9,0.9,0.9}{
            \makecell[{{p{\linewidth}}}]{
                \texttt{\tiny{[GM$|$GM]}}
                {'word': 'galva', 'assignment': 'team'}
            }
        }
    }
    & & \\ \\

    \theutterance \stepcounter{utterance}  
    & & \multicolumn{4}{p{0.6\linewidth}}{
        \cellcolor[rgb]{0.9,0.9,0.9}{
            \makecell[{{p{\linewidth}}}]{
                \texttt{\tiny{[GM$\rangle$P2]}}
                \texttt{Vārds saturns tika uzminēts pareizi. Vārds telpa tika uzminēts, bet tas ir nevainojams vārds. Jūsu gājiens beidzās tur.} \\
\texttt{Tagad sniedziet vēl vienu ar komatiem atdalītu sarakstu, kurā ir vismaz 1 un ne vairāk kā 2 vārdi no turpmāk minētā vārdu saraksta, kas vislabāk sasaucas vai ir visciešāk saistīti ar vārdu 'zilonis'. Neaizmirstiet atbildi sākt ar MINĒJUMS: . Piezīme: salīdzinājumā ar iepriekšējiem pieprasījumiem daži vārdi no saraksta ir izņemti.} \\
\\ 
\texttt{lāpa, mārciņa, pārsegs, vēstniecība, ziloņkauls, kļava, bērns, savienojums, zāle, ceļojums, kaķis, pātagu, sniegs, nāve, automašīna, galva, nu} \\
            }
        }
    }
    & \\ \\

    \theutterance \stepcounter{utterance}  
    & & & \multicolumn{4}{p{0.6\linewidth}}{
        \cellcolor[rgb]{0.9,0.9,0.9}{
            \makecell[{{p{\linewidth}}}]{
                \texttt{\tiny{[GM$\langle$P2]}}
                \texttt{MINĒJUMS: ziloņkauls, galva} \\
            }
        }
    }
     \\ \\

    \theutterance \stepcounter{utterance}  
    & & & \multicolumn{2}{p{0.3\linewidth}}{
        \cellcolor[rgb]{0.9,0.9,0.9}{
            \makecell[{{p{\linewidth}}}]{
                \texttt{\tiny{[GM$|$GM]}}
                ['ziloņkauls', 'galva']
            }
        }
    }
    & & \\ \\

    \theutterance \stepcounter{utterance}  
    & & & \multicolumn{2}{p{0.3\linewidth}}{
        \cellcolor[rgb]{0.9,0.9,0.9}{
            \makecell[{{p{\linewidth}}}]{
                \texttt{\tiny{[GM$|$GM]}}
                {'word': 'ziloņkauls', 'assignment': 'team'}
            }
        }
    }
    & & \\ \\

    \theutterance \stepcounter{utterance}  
    & & & \multicolumn{2}{p{0.3\linewidth}}{
        \cellcolor[rgb]{0.9,0.9,0.9}{
            \makecell[{{p{\linewidth}}}]{
                \texttt{\tiny{[GM$|$GM]}}
                {'word': 'ziloņkauls', 'assignment': 'team'}
            }
        }
    }
    & & \\ \\

    \theutterance \stepcounter{utterance}  
    & & & \multicolumn{2}{p{0.3\linewidth}}{
        \cellcolor[rgb]{0.9,0.9,0.9}{
            \makecell[{{p{\linewidth}}}]{
                \texttt{\tiny{[GM$|$GM]}}
                {'word': 'galva', 'assignment': 'team'}
            }
        }
    }
    & & \\ \\

    \theutterance \stepcounter{utterance}  
    & & & \multicolumn{2}{p{0.3\linewidth}}{
        \cellcolor[rgb]{0.9,0.9,0.9}{
            \makecell[{{p{\linewidth}}}]{
                \texttt{\tiny{[GM$|$GM]}}
                {'word': 'galva', 'assignment': 'team'}
            }
        }
    }
    & & \\ \\

    \theutterance \stepcounter{utterance}  
    & & & \multicolumn{2}{p{0.3\linewidth}}{
        \cellcolor[rgb]{0.9,0.9,0.9}{
            \makecell[{{p{\linewidth}}}]{
                \texttt{\tiny{[GM$|$GM]}}
                {'word': 'mārciņa', 'assignment': 'opponent'}
            }
        }
    }
    & & \\ \\

    \theutterance \stepcounter{utterance}  
    & & & \multicolumn{2}{p{0.3\linewidth}}{
        \cellcolor[rgb]{0.9,0.9,0.9}{
            \makecell[{{p{\linewidth}}}]{
                \texttt{\tiny{[GM$|$GM]}}
                {'word': 'vēstniecība', 'assignment': 'opponent'}
            }
        }
    }
    & & \\ \\

    \theutterance \stepcounter{utterance}  
    & & & \multicolumn{2}{p{0.3\linewidth}}{
        \cellcolor[rgb]{0.9,0.9,0.9}{
            \makecell[{{p{\linewidth}}}]{
                \texttt{\tiny{[GM$|$GM]}}
                {'hidden': {'team': ['nu', 'pārsegs', 'savienojums', 'zāle', 'pātagu'], 'innocent': ['ceļojums', 'bērns', 'kaķis', 'automašīna', 'lāpa'], 'opponent': ['sniegs', 'kļava'], 'assassin': ['nāve']}, 'revealed': {'team': {'team': ['atlantīda', 'saturns', 'ziloņkauls', 'galva'], 'innocent': ['amazone', 'telpa'], 'opponent': [], 'assassin': []}, 'opponent': {'team': [], 'innocent': [], 'opponent': ['pudele', 'dēlis', 'antarktīda', 'aste', 'mārciņa', 'vēstniecība'], 'assassin': []}}}
            }
        }
    }
    & & \\ \\

    \theutterance \stepcounter{utterance}  
    & & \multicolumn{4}{p{0.6\linewidth}}{
        \cellcolor[rgb]{0.9,0.9,0.9}{
            \makecell[{{p{\linewidth}}}]{
                \texttt{\tiny{[P1$\langle$GM]}}
                \texttt{Vārdi ziloņkauls, galva tika uzminēti pareizi. Jūsu komandas biedra gājiens beidzās tur.} \\
\texttt{Tagad sniedziet vēl vienu norādi, kas attiecas uz dažiem no jūsu atlikušajiem komandas vārdiem, un saistīto mērķa vārdu sarakstu. Neaizmirstiet sākt savu norādi ar 'PADOMS: ', pēc tam ievietot jaunu rindu un sākt ar komatiem atdalīto mērķa vārdu sarakstu ar 'MĒRĶI: '. Piezīme: salīdzinājumā ar iepriekšējiem pieprasījumiem daži vārdi no sarakstiem ir izņemti.} \\
\\ 
\texttt{Jūsu atlikušie komandas vārdi ir: nu, pārsegs, savienojums, zāle, pātagu.} \\
\texttt{Atlikušie pretinieka vārdi ir: sniegs, kļava, vēstniecība, mārciņa.} \\
\texttt{Atlikušie traucētājvārdi ir: ceļojums, bērns, kaķis, automašīna, lāpa.} \\
\texttt{Atlikušie slepkavības vārdi ir: nāve.} \\
            }
        }
    }
    & \\ \\

    \theutterance \stepcounter{utterance}  
    & \multicolumn{4}{p{0.6\linewidth}}{
        \cellcolor[rgb]{0.9,0.9,0.9}{
            \makecell[{{p{\linewidth}}}]{
                \texttt{\tiny{[P1$\rangle$GM]}}
                \texttt{PADOMS: pļava} \\
\texttt{MĒRĶI: zāle, pātagu} \\
            }
        }
    }
    & & \\ \\

    \theutterance \stepcounter{utterance}  
    & & & \multicolumn{2}{p{0.3\linewidth}}{
        \cellcolor[rgb]{0.9,0.9,0.9}{
            \makecell[{{p{\linewidth}}}]{
                \texttt{\tiny{[GM$|$GM]}}
                \texttt{pļava} \\
            }
        }
    }
    & & \\ \\

    \theutterance \stepcounter{utterance}  
    & & & \multicolumn{2}{p{0.3\linewidth}}{
        \cellcolor[rgb]{0.9,0.9,0.9}{
            \makecell[{{p{\linewidth}}}]{
                \texttt{\tiny{[GM$|$GM]}}
                ['zāle', 'pātagu']
            }
        }
    }
    & & \\ \\

    \theutterance \stepcounter{utterance}  
    & & & \multicolumn{2}{p{0.3\linewidth}}{
        \cellcolor[rgb]{0.9,0.9,0.9}{
            \makecell[{{p{\linewidth}}}]{
                \texttt{\tiny{[GM$|$GM]}}
                {'word': 'zāle', 'assignment': 'team'}
            }
        }
    }
    & & \\ \\

    \theutterance \stepcounter{utterance}  
    & & & \multicolumn{2}{p{0.3\linewidth}}{
        \cellcolor[rgb]{0.9,0.9,0.9}{
            \makecell[{{p{\linewidth}}}]{
                \texttt{\tiny{[GM$|$GM]}}
                {'word': 'pātagu', 'assignment': 'team'}
            }
        }
    }
    & & \\ \\

    \theutterance \stepcounter{utterance}  
    & & \multicolumn{4}{p{0.6\linewidth}}{
        \cellcolor[rgb]{0.9,0.9,0.9}{
            \makecell[{{p{\linewidth}}}]{
                \texttt{\tiny{[GM$\rangle$P2]}}
                \texttt{Vārdi ziloņkauls, galva tika uzminēti pareizi. Jūsu gājiens beidzās tur.} \\
\texttt{Tagad sniedziet vēl vienu ar komatiem atdalītu sarakstu, kurā ir vismaz 1 un ne vairāk kā 2 vārdi no turpmāk minētā vārdu saraksta, kas vislabāk sasaucas vai ir visciešāk saistīti ar vārdu 'pļava'. Neaizmirstiet atbildi sākt ar MINĒJUMS: . Piezīme: salīdzinājumā ar iepriekšējiem pieprasījumiem daži vārdi no saraksta ir izņemti.} \\
\\ 
\texttt{lāpa, pārsegs, kļava, bērns, savienojums, zāle, ceļojums, kaķis, pātagu, sniegs, nāve, automašīna, nu} \\
            }
        }
    }
    & \\ \\

    \theutterance \stepcounter{utterance}  
    & & & \multicolumn{4}{p{0.6\linewidth}}{
        \cellcolor[rgb]{0.9,0.9,0.9}{
            \makecell[{{p{\linewidth}}}]{
                \texttt{\tiny{[GM$\langle$P2]}}
                \texttt{MINĒJUMS: zāle, kļava} \\
            }
        }
    }
     \\ \\

    \theutterance \stepcounter{utterance}  
    & & & \multicolumn{2}{p{0.3\linewidth}}{
        \cellcolor[rgb]{0.9,0.9,0.9}{
            \makecell[{{p{\linewidth}}}]{
                \texttt{\tiny{[GM$|$GM]}}
                ['zāle', 'kļava']
            }
        }
    }
    & & \\ \\

    \theutterance \stepcounter{utterance}  
    & & & \multicolumn{2}{p{0.3\linewidth}}{
        \cellcolor[rgb]{0.9,0.9,0.9}{
            \makecell[{{p{\linewidth}}}]{
                \texttt{\tiny{[GM$|$GM]}}
                {'word': 'zāle', 'assignment': 'team'}
            }
        }
    }
    & & \\ \\

    \theutterance \stepcounter{utterance}  
    & & & \multicolumn{2}{p{0.3\linewidth}}{
        \cellcolor[rgb]{0.9,0.9,0.9}{
            \makecell[{{p{\linewidth}}}]{
                \texttt{\tiny{[GM$|$GM]}}
                {'word': 'zāle', 'assignment': 'team'}
            }
        }
    }
    & & \\ \\

    \theutterance \stepcounter{utterance}  
    & & & \multicolumn{2}{p{0.3\linewidth}}{
        \cellcolor[rgb]{0.9,0.9,0.9}{
            \makecell[{{p{\linewidth}}}]{
                \texttt{\tiny{[GM$|$GM]}}
                {'word': 'kļava', 'assignment': 'opponent'}
            }
        }
    }
    & & \\ \\

    \theutterance \stepcounter{utterance}  
    & & & \multicolumn{2}{p{0.3\linewidth}}{
        \cellcolor[rgb]{0.9,0.9,0.9}{
            \makecell[{{p{\linewidth}}}]{
                \texttt{\tiny{[GM$|$GM]}}
                \texttt{kļava} \\
            }
        }
    }
    & & \\ \\

    \theutterance \stepcounter{utterance}  
    & & & \multicolumn{2}{p{0.3\linewidth}}{
        \cellcolor[rgb]{0.9,0.9,0.9}{
            \makecell[{{p{\linewidth}}}]{
                \texttt{\tiny{[GM$|$GM]}}
                {'word': 'sniegs', 'assignment': 'opponent'}
            }
        }
    }
    & & \\ \\

    \theutterance \stepcounter{utterance}  
    & & & \multicolumn{2}{p{0.3\linewidth}}{
        \cellcolor[rgb]{0.9,0.9,0.9}{
            \makecell[{{p{\linewidth}}}]{
                \texttt{\tiny{[GM$|$GM]}}
                {'hidden': {'team': ['nu', 'pārsegs', 'savienojums', 'pātagu'], 'innocent': ['ceļojums', 'bērns', 'kaķis', 'automašīna', 'lāpa'], 'opponent': [], 'assassin': ['nāve']}, 'revealed': {'team': {'team': ['atlantīda', 'saturns', 'ziloņkauls', 'galva', 'zāle'], 'innocent': ['amazone', 'telpa'], 'opponent': ['kļava'], 'assassin': []}, 'opponent': {'team': [], 'innocent': [], 'opponent': ['pudele', 'dēlis', 'antarktīda', 'aste', 'mārciņa', 'vēstniecība', 'sniegs'], 'assassin': []}}}
            }
        }
    }
    & & \\ \\

    \theutterance \stepcounter{utterance}  
    & & \multicolumn{4}{p{0.6\linewidth}}{
        \cellcolor[rgb]{0.9,0.9,0.9}{
            \makecell[{{p{\linewidth}}}]{
                \texttt{\tiny{[P1$\langle$GM]}}
                \texttt{Vārds zāle tika uzminēts pareizi. Vārds kļava tika uzminēts, bet tas ir pretinieks vārds. Jūsu komandas biedra gājiens beidzās tur.} \\
\texttt{Tagad sniedziet vēl vienu norādi, kas attiecas uz dažiem no jūsu atlikušajiem komandas vārdiem, un saistīto mērķa vārdu sarakstu. Neaizmirstiet sākt savu norādi ar 'PADOMS: ', pēc tam ievietot jaunu rindu un sākt ar komatiem atdalīto mērķa vārdu sarakstu ar 'MĒRĶI: '. Piezīme: salīdzinājumā ar iepriekšējiem pieprasījumiem daži vārdi no sarakstiem ir izņemti.} \\
\\ 
\texttt{Jūsu atlikušie komandas vārdi ir: nu, pārsegs, savienojums, pātagu.} \\
\texttt{Atlikušie pretinieka vārdi ir: sniegs.} \\
\texttt{Atlikušie traucētājvārdi ir: ceļojums, bērns, kaķis, automašīna, lāpa.} \\
\texttt{Atlikušie slepkavības vārdi ir: nāve.} \\
            }
        }
    }
    & \\ \\

    \theutterance \stepcounter{utterance}  
    & \multicolumn{4}{p{0.6\linewidth}}{
        \cellcolor[rgb]{0.9,0.9,0.9}{
            \makecell[{{p{\linewidth}}}]{
                \texttt{\tiny{[P1$\rangle$GM]}}
                \texttt{PADOMS: vads} \\
\texttt{MĒRĶI: savienojums, pātagu} \\
            }
        }
    }
    & & \\ \\

    \theutterance \stepcounter{utterance}  
    & & & \multicolumn{2}{p{0.3\linewidth}}{
        \cellcolor[rgb]{0.9,0.9,0.9}{
            \makecell[{{p{\linewidth}}}]{
                \texttt{\tiny{[GM$|$GM]}}
                \texttt{vads} \\
            }
        }
    }
    & & \\ \\

    \theutterance \stepcounter{utterance}  
    & & & \multicolumn{2}{p{0.3\linewidth}}{
        \cellcolor[rgb]{0.9,0.9,0.9}{
            \makecell[{{p{\linewidth}}}]{
                \texttt{\tiny{[GM$|$GM]}}
                ['savienojums', 'pātagu']
            }
        }
    }
    & & \\ \\

    \theutterance \stepcounter{utterance}  
    & & & \multicolumn{2}{p{0.3\linewidth}}{
        \cellcolor[rgb]{0.9,0.9,0.9}{
            \makecell[{{p{\linewidth}}}]{
                \texttt{\tiny{[GM$|$GM]}}
                {'word': 'savienojums', 'assignment': 'team'}
            }
        }
    }
    & & \\ \\

    \theutterance \stepcounter{utterance}  
    & & & \multicolumn{2}{p{0.3\linewidth}}{
        \cellcolor[rgb]{0.9,0.9,0.9}{
            \makecell[{{p{\linewidth}}}]{
                \texttt{\tiny{[GM$|$GM]}}
                {'word': 'pātagu', 'assignment': 'team'}
            }
        }
    }
    & & \\ \\

    \theutterance \stepcounter{utterance}  
    & & & \multicolumn{2}{p{0.3\linewidth}}{
        \cellcolor[rgb]{0.9,0.9,0.9}{
            \makecell[{{p{\linewidth}}}]{
                \texttt{\tiny{[GM$|$GM]}}
                \texttt{opponent has won} \\
            }
        }
    }
    & & \\ \\

\end{supertabular}
}

\end{document}
